\section{Security Analysis}
\label{sec:security}

\subsection{Threat Model}

The adversary $\mathcal{A}$ is computationally bounded and may: eavesdrop, replay,
delay, inject, and modify V2V channel messages; deploy rogue RSUs or impersonate the
AIS; compromise individual vehicles (but not the TA); attempt Sybil registrations;
and observe AST request timing.

\subsection{Authentication Soundness}

\begin{theorem}[Authentication Soundness]
Under the $q$-Strong Bilinear Diffie-Hellman ($q$-SBDH) assumption over BN254, no
computationally bounded adversary can produce a valid Groth16 proof $\pi$ with
respect to Merkle root $R_{current}$ without possessing a valid $\mathsf{AST}_i$
whose leaf $\mathsf{leaf}_i$ is present in the Merkle tree with root $R_{current}$.
\end{theorem}
\begin{proof}[Sketch]
Groth16 is a proof of knowledge under the $q$-SBDH assumption~\cite{groth2016}.
The knowledge extractor guarantees that a valid proof $\pi$ for public inputs
$(R, t_{current}, h_m)$ implies existence of a witness $(\mathsf{AST}_i, \pi^{MT}_i,
m)$ satisfying all three circuit constraints. Forging without this witness requires
breaking $q$-SBDH in $\mathbb{G}_1$.
\end{proof}

\subsection{AIS Authenticity and Rogue Infrastructure Prevention}

\begin{theorem}[AIS Authenticity]
If ECDSA is existentially unforgeable under chosen-message attack (EUF-CMA), no
adversary without $sk_{TA}$ can produce a valid $\mathsf{cert}_{AIS}$ accepted by
an honest vehicle, and no adversary without $sk_{AIS}$ can produce a valid
$\mathsf{sig}_{AIS}$ accepted during AST acquisition.
\end{theorem}
\begin{proof}[Sketch]
$\mathsf{cert}_{AIS}$ is an ECDSA signature under $sk_{TA}$; forging it reduces
to breaking ECDSA by EUF-CMA. A rogue AIS lacking $sk_{TA}$ is rejected at
Phase~2 Step~1. Similarly, $\mathsf{sig}_{AIS}$ requires $sk_{AIS}$ to forge.
\end{proof}

\subsection{Privacy: Unlinkability of Safety Messages}

\begin{theorem}[Unlinkability]
The zero-knowledge property of Groth16 guarantees that proofs $\pi_1, \pi_2$ from
the same vehicle (same $\mathsf{AST}_i$, fresh randomness) are computationally
indistinguishable from proofs generated by two independent vehicles.
\end{theorem}
\begin{proof}[Sketch]
The Groth16 zero-knowledge simulator produces proofs indistinguishable from honest
proofs without the witness. Independent randomisation of each proof ensures the
joint distribution of $(\pi_1, \pi_2)$ reveals no more than two independently
simulated proofs. The public inputs $(R, t_{current}, h_m)$ reveal only that the
prover holds \emph{some} valid AST in the epoch (shared by all valid vehicles);
$sid_i$ and $\pi^{MT}_i$ remain in the private witness.
\end{proof}

\noindent\textit{Note on AST-level linkability:} Multiple proofs under the same AST
share the same public $R$. Since $R$ is shared by all valid vehicles in the epoch,
it does not identify the prover; fresh proof randomness ensures consecutive proofs
from the same vehicle are indistinguishable from those of different vehicles.

\subsection{Replay Attack Prevention}

\begin{theorem}[Replay Resistance]
A proof $\pi$ generated for message $m$ at timestamp $t_1$ cannot be validly
replayed for a different timestamp $t_2 \neq t_1$ or message $m' \neq m$.
\end{theorem}
\begin{proof}[Sketch]
$h_m = H_{Poseidon}(m \| t_{current})$ and $t_{current}$ are public inputs; circuit
constraint~(3) enforces $h_m = H_{Poseidon}(m \| t_{current})$ in-circuit. Changing
$t_2$ or $m'$ while reusing $\pi$ creates a public-input mismatch. Generating a
fresh valid proof for $(m', t_2)$ without $\mathsf{AST}_i$ is infeasible by
Theorem~1. Verifiers additionally reject proofs with
$|t_{verifier} - t_{current}| > \Delta_t$ ($\Delta_t \approx 200$\,ms for clock
drift tolerance).
\end{proof}

\subsection{Revocation Correctness}

\begin{theorem}[Revocation Correctness]
After revocation of $\mathsf{AST}_i$, any proof generated by $V_i$ referencing the
updated Merkle root $R_{new}$ will fail verification at all honest verifiers.
\end{theorem}
\begin{proof}[Sketch]
Revocation removes $\mathsf{leaf}_i$ and recomputes $R_{new}$. Any Merkle path
$\pi^{MT}_i$ for the old leaf does not hash to $R_{new}$, causing circuit
constraint~(1) to fail. The verifier's root-matching check rejects before batch
verification. During $T_{grace}$, proofs referencing $R_{prev}$ are accepted only
if the leaf was present in $R_{prev}$; revocation excludes the leaf from $R_{new}$
at the next epoch, fully excluding the vehicle once $T_{grace}$ expires.
\end{proof}

\subsection{Sybil Resistance}

Sybil resistance derives from the TA's vehicle registration process (Phase~1),
which verifies real-world identity (VIN) before issuing $\sigma_i$. Multiple
registrations require distinct legitimate VINs, bounded by the physical vehicle
supply. This matches the trust assumption of SCMS and introduces no additional
assumptions beyond the existing PKI model.

\subsection{Timing Fingerprinting Resistance}

The jitter mechanism ($\delta \sim \mathcal{U}(0, 30\text{s})$ for non-urgent AST
requests) provides 30\,s of timing ambiguity, preventing precise timing correlation
of AST request patterns to vehicle identities at the cost of at most 30\,s delay
in non-urgent requests.

\subsection{Emergency Message Unlinkability}
\label{sec:emergency_security}

A natural question is whether emergency BSMs (emergency braking, forward collision
warning) require a fallback to long-term ECDSA credentials for low latency. We
argue they must not. If vehicle $V_i$ uses a fixed long-term key $pk_i$ for any
message class, an observer collecting emergency BSMs over time can link all
emergency events to the same vehicle by the common $pk_i$---constructing a complete
emergency history even without knowing the real-world identity behind $pk_i$.
In SCMS, pseudonym certificates partially mitigate this within a rotation period,
but a vehicle with frequent emergency events remains linkable across events
within the same certificate window.

\textbf{Emergency slot pool.}
ULP-V2V-Auth addresses this by applying ZKP authentication uniformly to all
message classes. A dedicated cache of $N_e$ proof slots is pre-committed to
emergency-class BSM templates (emergency flag, last-known position at 1\,s
resolution). Upon an emergency event:
\begin{enumerate}
    \item Dequeue next emergency slot from pool: $< 0.01$\,ms
    \item Bind $h_m = H_{\text{Poseidon}}(m_{\text{emergency}} \| t_{\text{current}})$: $0.437$\,ms
    \item Broadcast immediately: $0$\,ms
\end{enumerate}
Total emergency authentication latency: $\mathbf{0.439}$\,ms, identical to routine
BSMs. On production hardware (NXP~S32G, $\leq$98\,ms/slot), proofs are generated
on-demand within one BSM interval, removing the pre-commitment requirement entirely.

\begin{theorem}[Emergency Message Unlinkability]
Emergency BSMs authenticated via the ZKP proof-slot model are computationally
unlinkable from routine BSMs and from each other, under the zero-knowledge
property of Groth16~\cite{groth2016}.
\end{theorem}
\begin{proof}[Sketch]
Each emergency proof uses fresh, independently sampled Groth16 randomness. The
ZK simulator $\mathcal{S}$ produces transcripts computationally indistinguishable
from honest proofs without knowledge of the witness $(\mathsf{AST}_i,
m_{\text{emergency}}, \pi^{MT}_i)$. The emergency flag is a private witness field
element; the public input $h_m = H_{\text{Poseidon}}(m \| t)$ binds only the
message content and timestamp. No component of the public transcript identifies
the AST, session, or vehicle. An observer learns that \emph{some} valid vehicle
triggered an emergency at location $L$---no more information than is inherent in
the physical event itself and observable by any road user.
\end{proof}

\subsection{Batch Verification Fallback: DoS Analysis and DCV Mitigation}
\label{sec:dcv_security}

A well-known vulnerability of batch verification schemes is the \emph{batch
injection attack}: an adversary broadcasts one invalid proof among $k$ legitimate
proofs, causing the batch check to fail and forcing fallback verification. Under a
naive $O(k)$ individual fallback, this transforms the $4.65\times$ speedup into an
effective $4.65\times$ slowdown---a near-$22\times$ swing relative to the honest
case.

\textbf{Pre-screening.}
Before batching, every incoming proof $(A_j, B_j, C_j)$ is checked for: (i)~valid
$\mathbb{G}_1$ membership of $A_j, C_j$; (ii)~valid $\mathbb{G}_2$ membership of
$B_j$; (iii)~timestamp freshness $|t_{\mathit{verifier}} - t_j| \leq \Delta_t$;
and (iv)~correct Merkle root $R_j = R_{\mathit{current}}$. These checks cost
$< 0.1$\,ms per proof and reject trivially malformed injections at zero pairing
cost.

\textbf{DCV fallback bound.}
For injections that pass pre-screening (i.e., structurally valid but semantically
invalid proofs), Algorithm~\ref{alg:dcv_fallback} limits fallback cost to
$O(w \log k)$ sub-batch verifications~\cite{ferrara2008finding}. An adversary
injecting $w$ invalid proofs triggers at most $2w\lceil\log_2 k\rceil$ sub-batch
operations rather than $k$ individual verifications. At $k=30$, $w=1$: DCV costs
$\approx 500$\,ms versus $2{,}100$\,ms for naive fallback. The adversary's DoS
gain is bounded at $O(\log k)$ rather than $O(k)$.

\textbf{Rate limiting.}
IEEE 802.11p EDCA congestion control and the MAC-layer 10\,Hz BSM rate limit bound
the injection rate to at most one invalid proof per 100\,ms window per adversarial
vehicle. Combined with DCV, this caps the sustained CPU overhead of batch injection
at $\approx 500$\,ms per window---equal to one honest verification window under
moderate traffic ($k=30$).

This defence is not unique to ULP-V2V-Auth: the batch injection vulnerability and
DCV mitigation apply equally to certificateless batch signature
schemes~\cite{zhang2024}. We note it explicitly because no prior vehicular
authentication proposal has analysed or mitigated this attack in the ZKP setting.

Table~\ref{tab:security_summary} summarises the security properties established
in this section.

\begin{table}[t]
\centering
\caption{Security Properties of ULP-V2V-Auth}
\label{tab:security_summary}
\resizebox{\columnwidth}{!}{%
\begin{tabular}{@{}lll@{}}
\toprule
\textbf{Property} & \textbf{Guarantee} & \textbf{Assumption} \\
\midrule
Authentication soundness  & Forging proof requires breaking $q$-SBDH & $q$-SBDH on BN254 \\
AIS authenticity          & Rogue AIS rejected without $sk_{TA}$      & ECDSA EUF-CMA \\
Message unlinkability     & Proofs reveal no vehicle identity          & Groth16 ZK \\
Replay resistance         & $(m, t)$ binding via Poseidon in-circuit   & Collision resistance \\
Emergency unlinkability   & Emergency BSMs unlinkable from all others  & Groth16 ZK \\
Revocation correctness    & Revoked AST fails after $T_{grace}$       & Merkle binding \\
Sybil resistance          & Bounded by physical vehicle registry       & TA trust \\
Batch injection DoS       & DCV bounds fallback at $O(w \log k)$      & Pre-screening + DCV \\
\bottomrule
\end{tabular}%
}
\end{table}
